\documentclass{article}
\usepackage{amsfonts, amsmath, amsthm, amssymb, enumerate, fancyhdr, gensymb, lastpage, mathtools, parskip, graphicx}
\usepackage{xcolor, tikz-cd}
\usepackage[a4paper, total={6in, 8in}]{geometry}
\usepackage{setspace}

\usepackage[OT2,T1]{fontenc}
\DeclareSymbolFont{cyrletters}{OT2}{wncyr}{m}{n}
\DeclareMathSymbol{\Sha}{\mathalpha}{cyrletters}{"58}

% Theorem environments
\newtheorem{theorem}{Theorem}
\newtheorem{lemma}[theorem]{Lemma}
\newtheorem{proposition}[theorem]{Proposition}
\newtheorem{corollary}[theorem]{Corollary}

% Definition-style environments
\theoremstyle{definition}
\newtheorem{definition}[theorem]{Definition}
\newtheorem{example}[theorem]{Example}

% Remark-style environments
\theoremstyle{remark}
\newtheorem{remark}[theorem]{Remark}


%\onehalfspacing
\setstretch{1.3}
\newcommand{\wg}[1]{\textcolor{violet}{#1}}

%\renewcommand{\C}{\mathcal C}
\newcommand{\Z}{\mathbb Z}
\newcommand{\Ztwo}{\mathbb Z / 2\mathbb Z}
\newcommand{\Ztwosq}{(\mathbb Z / 2 \mathbb Z)^2}
\newcommand{\GalK}{G_{\bar K / K}}
\newcommand{\Q}{\mathbb Q}
\newcommand{\R}{\mathbb R}
\newcommand{\im}{\text{im }}
\renewcommand{\div}{\, \vert \,}
\newcommand{\F}{\mathbb F}
\newcommand{\EQW}{E'(\Q)/2E(\Q)}
\newcommand{\Etors}{E(\Q)_{\text{tors}}}
\renewcommand{\F}[1]{\mathbb F_{#1}}
\newcommand{\Qst}{\mathbb Q(S,2)}
\newcommand{\Sphi}{S^\phi}
\newcommand{\Sphihat}{S^{\hat \phi}}
\newcommand{\bark}{{\bar{k}}}
\newcommand{\Qsqs}{(\Q^\times)/(\Q^\times)^2}
\newcommand{\Zn}[1]{\mathbb Z / #1 \mathbb Z}

\newenvironment{wgenv}{\color{violet}}{}

\newcommand{\smalltwobytwo}[4]{
\left( \begin{smallmatrix*}
   #1 & #2 \\ #3 & #4 
\end{smallmatrix*}\right) 
}
\newcommand{\bmat}[1]{
    \begin{bmatrix*}
        #1
    \end{bmatrix*}
}

\title{
    Composition Series
}
\author{Will Gilroy}
\date{July 2026}

\begin{document}
\maketitle
\wg{As of $18th$ July this note is still in rough form}

Let $R$ be a local ring with maximal ideal $\mathfrak m = (y_1, \cdots, y_m)$ and let $M =
R/I$ be an $R$-module where $I$ is $\mathfrak m$-primary. We want to try and
construct a composition series for $M$ by iteratively constructing a chain where the partial
quotient is $R/I'$ where $I' = (\text{gens } I, y_1)$. Let's try to work out how
this might work with an example

Let $R = k[x]_{(x)}$ and $M = R/(x^3)$. 
The correspondence theorem gives us that the submodules of $M$ correspond to the
ideals $0, (x^2), (x^3)$, giving the submodules $0/(x^3), (x^2)/(x^3)$, and
$(x^3)/(x^3)$. 

\wg{Mike was giving us a procedure so that the first partial quotient would 
be something like $R/(x^3, x) = R/(x)$ which is now modding out by the maximal 
ideal and the procedure would terminate. I think we would then get the wrong
length by this procedure, but I admit I do not totally recall the 
initial chain that we would start with.}
\wg{Perhaps we can try this example again with $R = k[x,y]_{(x,y)}$ and $R/(x^2,
y^2)$?}

\appendix
\subsection{Correspondence Theorem and Submodules of $R/I$.}
A quick comment on how the correspondence theorem gives us the 
submodules of $R/(x^3)$. Recall the correspondence theorem for groups:
if $G$ is a group and $N \leq G$ a normal subgroup then there is a bijection
\[
    \{ \text{Subgroups of $G$ containing $N$} \} \leftrightarrow
    \{ \text{Subgroups of $G/N$}\}.
\]
This correspondence theorem extends to $R$ modules in the following way.
Let $R$ be a ring, viewed as an $R$ module, and let $N$ an submodule of $R$. 
Then the quotient $R$-module $R/N$ is in particular a quotient group of $R$.
Applying the correspondence theorem for groups gives us that the subgroups (and
thus the only contenders for the submodules of $R/N$)
of $R/N$ are $\{R/L : N \subset L \leq R \}$. We can extend all of these
subgroups to $R$ submodules of $R/N$ in the usual way and we thus have the
following bijection
\[
    \{\text{Submodules of $R$ containing $N$}\} \leftrightarrow
    \{\text{Submodules of $R/N$}\}. 
\]

In particular, if $M = R/I$ where $I \leq R$ is an ideal of $R$ applying this
bijection gives
\[
    \{\text{Ideals of $R$ containing $I$}\} \leftrightarrow
    \{\text{Submodules of $R/I$}\},
\]
noting that the submodules of $R$ are exactly its ideals.

\begin{example}
Let's apply this to the example in the main body above, let $R = k[x]$ and $M =
k[x]/(x^3)$. To construct a chain of $M$ we need to understand its submodules.
The submodules of $M$ are in bijection with the ideals of $R$ containing
$(x^3)$. These ideals are exactly $(1), (x), (x^2), (x^3)$. Applying the
bijection map \wg{(which we did not describe)} gives the submodules
\[
    R/(x^3),(x)/(x^3), (x^2)/(x^3), (x^3)/(x^3),
\]
where each of these quotients are the quotient modules 
(note for example that the ideal $(x)$ is a subgroup of $(x^3)$ when viewed
purely as additive groups.)
we can then construct a composition series
\begin{align*}
    (x^3)/(x^3) = 0 \subseteq (x^2)/(x^3) \subseteq (x)/(x^3) \subseteq R/(x^3) = (1)/(x^3).
\end{align*}
Since we have enumerated all of the submodules of $R/(x^3)$ we can see that
there is no longer chain which can be constructed (in particular, we have a
chain which has used all of the submodules) and so the above chain is in fact a
composition series (one can also check that the partial quotients are simple if
one prefers). Thus we have shown that the length 
\[
    \ell_{R}(R/(x^3)) = 3.
\]
Note that following a similar argument shows that $\ell_{R}(R/(x^n)) = n$. 
\wg{I think we can also further note that, since localisation is an exact
functor, the above chain remains a chain and so the length should not change 
under localising at $(x)$. We should work out the details for this.}
\end{example}


\end{document}